% Mục A.2. Phụ trách: TV1.
\phutrach{TV1 -- mục A.2; mỗi mục kết thúc bằng ảnh hưởng tới chỉ mục}

\subsection{Tách tầng tính toán và tầng lưu trữ}
\ghichu{TiDB server, TiKV, Coprocessor. Giải thích \code{TableReader}, \code{IndexReader} (covering), \code{IndexLookUp}, \code{Point\_Get}, vì Chương~\ref{sec:doc} đo chúng trực tiếp.}

\hinhcho{Sơ đồ 1 PD, 1 TiDB server, 3 TiKV đúng với cụm đã cài}{Kiến trúc cụm TiDB dùng trong đề tài}{fig:kien-truc}

\subsection{Ánh xạ bảng và chỉ mục sang Key-Value}

\subsubsection{Dòng dữ liệu và chỉ mục}
\ghichu{Dạng key của dòng, chỉ mục unique và không unique; đối chiếu với B+Tree phụ của InnoDB.}

\subsubsection{RocksDB và LSM-tree}
\ghichu{WAL, MemTable, SST, compaction. Nêu cả lợi ích (không sửa trang tại chỗ) và cái giá (khuếch đại ghi). Không kết luận bên nào rẻ hơn; RQ4 phân xử.}

\subsection{Region và Multi-Raft}
\ghichu{Kích thước Region theo \code{SHOW CONFIG} trên v8.5.5. Quorum là đa số peer của một Region. Định nghĩa C3 và C1 dùng Bảng~\ref{tab:ban-sao}. Ba TiKV là ba container trên một máy. Bỏ số ``3--8 giây'' ở bản cũ.}

\subsection{Placement Driver, TSO và giao dịch Percolator}
\ghichu{TSO, prewrite, commit, MVCC. Mức cô lập Snapshot Isolation, không hỗ trợ SERIALIZABLE. Số key phải ghi trong một giao dịch tăng theo số chỉ mục.}

\subsection{Kết luận chương}
\begin{table}[!htp]
\centering\bangfont
\caption{Cơ chế của TiDB và câu hỏi nghiên cứu liên quan}\label{tab:co-che}
\begin{tabular}{|L{4.6cm}|L{6.6cm}|C{3.0cm}|}
\hline
\dongdau \textbf{Cơ chế} & \textbf{Ảnh hưởng tới chỉ mục} & \textbf{Câu hỏi} \\ \hline
IndexReader, IndexLookUp & \chuaco & RQ1, RQ2 \\ \hline
Ánh xạ KV, LSM-tree & \chuaco & B.1, RQ4 \\ \hline
Region, Multi-Raft & \chuaco & RQ4 (C1, C3) \\ \hline
TSO, Percolator & \chuaco & RQ1, RQ4 \\ \hline
\end{tabular}
\end{table}
